\chapter{Theoretical Background}
\label{ch:background}

% DRAFT BRIEF: protocol + methodology grounding so Ch.4-7 need not re-explain basics.
% Narrative + citations + one feature-matrix table; avoid code listings. No RTL dependency.
% Sources: phase1_spec_v2.md, MIPI I3C Basic v1.1.1, UVM 1.2, specs 02/03.

\section{\iic{} recap}
\brief{Electrical model (open-drain + pull-ups, wired-AND), START/STOP, 7-bit addressing
+ R/W, ACK/NACK, clock stretching. Frame the limitations I3C removes.}

\figph{Address byte}{9-bit address byte A6..A0 + RnW + ACK, MSB first}
  {phase1_spec_v2.md}{WaveDrom (bus/address_byte.json)}

\section{MIPI I3C Basic v1.1.1 (SDR, frame format)}
\brief{I3C electrical changes (\SI{1.8}{\volt} LV-CMOS push-pull for data), SDR mode, the
broadcast address 0x7E, the SDR private write/read frames, and T-bit semantics (odd-parity
bit in writes; end-of-data signalling bit in reads).}

\figph{SDR write frame}{S, addr+W, ACK, \{data+T\}$\times$N, P (T = odd parity)}
  {phase1_spec_v2.md}{WaveDrom (bus/sdr_write_frame.json)}
\figph{SDR read frame}{S, addr+R, ACK, \{data,T=1\}...\{data,T=0\}, P (T = end-of-data)}
  {phase1_spec_v2.md}{WaveDrom (bus/sdr_read_frame.json)}

\section{Bus conditions (START / STOP / Sr)}
\brief{Define START (SDA fall while SCL high), STOP (SDA rise while SCL high),
repeated-START, and bus-available/idle timing. Tie to \texttt{bus_monitor} and
\texttt{scl_generator} later.}

\figph{Bus conditions}{START / repeated-START / STOP --- SDA edge relative to SCL high}
  {phase1_spec_v2.md}{WaveDrom (bus/bus_conditions.json)}

\section{Open-drain vs push-pull}
\brief{Why address/ACK phases stay open-drain (legacy arbitration, hot devices) while data
phases go push-pull (speed). Define the OD-to-PP boundary and the role of a \texttt{sel_od_pp} select.}

\figph{Open-drain vs push-pull phases}
  {OD (addr/ACK) vs PP (data); OD-to-PP boundary; \texttt{sel_od_pp} signal}
  {phase1_spec_v2.md}{WaveDrom (bus/od_vs_pp.json)}

\section{Dynamic address assignment (ENTDAA)}
\brief{The ENTDAA algorithm from the master's perspective: broadcast 0x7E+W, the ENTDAA CCC,
Sr, then per-device 0x7E+R reading 64-bit \{PID, BCR, DCR\}, assigning a dynamic address +
parity, ACK, looping until no device responds. Sets up Ch.5.7's two-FSM implementation.}

\figph{ENTDAA arbitration}
  {7E+W, ACK, CCC, Sr, 7E+R, 64-bit PID+BCR+DCR, dyn-addr+parity, ACK loop}
  {phase1_spec_v2.md}{WaveDrom (bus/entdaa_arbitration.json)}

\section{CCC overview}
\brief{Common Command Codes: broadcast vs direct framing; the supported subset
(ENTDAA 0x07, ENEC/DISEC broadcast + direct).}

\figph{CCC frame}{broadcast + direct CCC frame format}
  {phase1_spec_v2.md}{WaveDrom (bus/ccc_frame.json)}

\section{I3C-vs-\iic{} comparison}
\brief{A consolidated feature-matrix table: speed, electrical, addressing, IBI, CCCs,
in-band interrupts, backward compatibility. Mark in-scope rows.}

\tabph{I3C-vs-\iic{} feature matrix}{MIPI I3C Basic v1.1.1}
\tabph{SDR + \iic{}-FM timing parameters}{phase1_spec_v2.md}

\section{UVM 1.2 methodology overview}
\brief{Component model (test, env, agent \{driver/monitor/sequencer\}, sequence), TLM
analysis ports, \texttt{uvm_config_db}, factory. Just enough to ground Ch.6--7. Cite
Accellera UVM 1.2.}

\figph{UVM testbench layering}
  {test to env to agents to sequencer/driver/monitor; TLM ports}{conceptual}{TikZ}
